\documentclass[10pt,a4paper]{amsart}
\usepackage[margin=19mm]{geometry}
\usepackage{amsmath,amssymb,mathtools}
\usepackage[scr=rsfs,cal=euler]{mathalfa}
\usepackage{xfrac}
\usepackage[unicode,hidelinks,bookmarks=false]{hyperref}
\newcommand{\ie}{i.e.\ }
\newcommand{\cf}{cf.\ }
\newcommand{\resp}{respectively\ }
\newcommand{\Subsection}{\subsection}
\providecommand{\nobreakdash}{\nobreak\hbox{-}}
\setlength{\emergencystretch}{2em}
\multlinegap0pt
\usepackage{algorithm}\usepackage{algpseudocode}
\newtheorem{lemma}{Lemma}
\newtheorem{proposition}{Proposition}
\title{Exact Volumes of Semi-Algebraic Convex Bodies}
\author{Lakshmi Ramesh, Nicolas Weiss}
\date{Source: arXiv:2602.04707v2 (CC BY 4.0)}
\begin{document}
\maketitle

\noindent\emph{Source and licence.} Exact Volumes of Semi-Algebraic Convex Bodies. Version arXiv:2602.04707v2 (CC BY 4.0), distributed by arXiv under the Creative Commons Attribution 4.0 International licence, http://creativecommons.org/licenses/by/4.0/, as recorded in the arXiv licence metadata for this exact version and shown on the abstract page https://arxiv.org/abs/2602.04707v2. Full licence notice: \texttt{licenses/arxiv-2602.04707-CC-BY-4.0.txt}.\par\smallskip

\noindent\textit{Editorial scope.} Counted unit: the article's proposition prop:deformed-intersection (source0.tex lines 442--446). The article's complete original proof of it is delivered with it (source0.tex lines 447--456, one \texttt{proof} environment of 10 lines). No mathematics is written, repaired or re-derived: the statement and the proof are the article's own lines, in the article's own notation, and no retained line is substituted or re-flowed.  Retained uncounted as the source-local context the proof needs: lines 239--241; lines 374--376; lines 379--385.  Nothing was omitted from any retained span, so the package carries no omission record.  No retained line contains a citation, so the wrapper carries no bibliography and no bibliography entry.  No retained line contains a figure environment, an image-inclusion command or drawing code, so no image is delivered and the package declares no asset dependency.  Editorial formatting only, in addition: an AMS \texttt{amsart} preamble that restates the article's own declarations and macros used by the retained lines, namely the environments \texttt{lemma}, \texttt{proposition} on separate counters, together with the standard packages those lines need.  The retained environments print in this document as Lemma 1, Proposition 1 in order of appearance, whereas the article prints the same items with the numbering of its own full document.  The complete native source of the article as distributed in the arXiv e-print for this exact version is bundled unchanged as \texttt{sources/193989/source0.tex}.
\begin{equation}\label{eq:deformed-intersection}
    C_t = \{x \in \mathbb R^n \mid \prod_i f_i(x) -t > 0 \} \cap C. 
\end{equation}

\begin{equation}\label{eqdef:concave}
    f\big(\alpha x + (1 - \alpha)y\big) \geq \alpha f(x) + (1-\alpha)f(y).
\end{equation}

\begin{lemma}\label{lemma:superlevelsets}
    Let $U \subset \mathbb{R}^n$ be a convex region and $f: U \rightarrow \mathbb{R}$ a concave function on that region. Then for any $t$, the super-levelset
    \begin{equation}
        f_{>t} := \{x \in \mathbb{R}^n \mid f(x) > t\}
    \end{equation}
    is a convex subset of $\mathbb{R}^n$.
\end{lemma}

\begin{proposition}\label{prop:deformed-intersection}
    Let $C$ be the common positivity locus of concave polynomials $f_1,\ldots,f_k \in \mathbb{Q}[\mathbf x]$. Then,
    its deformation $C_t$ as defined in \eqref{eq:deformed-intersection} is a convex set for all $t>0$.
    % 
\end{proposition}

\begin{proof}
    Recall that a positive function $f$ is defined to be log-concave if $\log(f)$ is also concave. Observe then that Lemma~\ref{lemma:superlevelsets} also holds for log-concave functions since the logarithm is a strictly monotonic function.
    
    In our case, the deformations $C_t$ are the super-levelsets of the product $F:=\prod^k_{i=1} f_i$ when restricted to $C$. But $C$ is the common positivity locus of all $f_i$, where automatically each of the $f_i$ is log-concave.
    Then, log-concavity of $F$ follows since 
    \begin{equation*}
        \log(F) = \log(f_1) + \ldots + \log(f_k)
    \end{equation*}
    and since sums of concave functions are concave by the defining Equation~\eqref{eqdef:concave}.
\end{proof}

\end{document}
